%
% Copyright (c) 2021-2023 Zeping Lee
% Released under the MIT License.
% Repository: https://github.com/zepinglee/citeproc-lua
%

\documentclass{l3doc}

\usepackage{mathpazo}
\usepackage{helvet}
\usepackage{listings}

\NewDocumentCommand\opt{m}{\texttt{#1}}

\lstdefinestyle{plain}{
  basicstyle = \ttfamily\small,
  basewidth  = 0.51em,
}

\lstnewenvironment{LaTeXdemo}{
  \lstset{
    style      = plain,
    frame      = single,
    gobble     = 2,
    language   = [LaTeX]TeX,
  }
}{}

\lstnewenvironment{bash}{
  \lstset{
    style      = plain,
    % gobble     = 2,
    language   = bash,
  }
}{}

% \hypersetup{hidelinks}
% \urlstyle{same}

\begin{document}

\title{%
  Bibliography formatting with \pkg{citation-style-language}
}

\author{%
  Zeping Lee%
  \thanks{%
    E-mail:
    \href{mailto:zepinglee@gmail.com}
      {zepinglee@gmail.com}%
  }%
}

\date{2023-04-02 v0.4.0}

\maketitle

% \begin{abstract}
%   Foo
% \end{abstract}

\begin{documentation}

\section{Introduction}

The Citation Style Language\footnote{\url{https://citationstyles.org/}} (CSL)
is an XML-based language that defines the formats of citations and bibliography.
There are currently thousands of styles in CSL including the most widely used
APA, Chicago, Vancouver, etc.
The \pkg{citation-style-language} package is aimed to provide another reference formatting method
for LaTeX that utilizes the CSL styles.
It contains a citation processor implemented in pure Lua (\pkg{citeproc-lua})
which reads bibliographic metadata and performs sorting and formatting on both
citations and bibliography according to the selected CSL style.
A LaTeX package (\file{citation-style-language.sty}) is provided to communicate with the processor.

Note that this project is in early development stage and some features of CSL
are not implemented yet.
Comments, suggestions, and bug reports are welcome.


\section{Installation}

This package is available from TeX Live 2022 or later versions.
For most users, the easiest way is to install it via |tlmgr|.
If you want to install the GitHub develop version of this package,
you may follow the steps below.

The \pkg{citation-style-language} requires the following packages:
\pkg{filehook}, \pkg{l3kernel}, \pkg{l3packages}, \pkg{lua-uca}, \pkg{lualibs},
\pkg{luatex}, \pkg{luaxml}, and \pkg{url}.
\pkg{l3build} is also required for actually performing the installation.
Make sure they are already installed in the TeX distribution.

\begin{bash}
  git clone https://github.com/zepinglee/citeproc-lua  # Clone the repository
  cd citeproc-lua
  git submodule update --init --remote                 # Fetch submodules
  l3build install
\end{bash}

These commands install the package files to |TEXMFHOME| which is usually
|~/texmf| on Linux or |~/Library/texmf| on macOS.
Besides, the |citeproc-lua| executable needs to be copied to some directory in
the |PATH| environmental variable so that it can be called directly in the shell.
For example provided |~/bin| is in |PATH|:

\begin{bash}
  cp citeproc/citeproc-lua.lua "~/bin/citeproc-lua"
\end{bash}

To uninstall the package from |TEXMFHOME|, just run |l3build uninstall|.

\section{Getting started}

An example of using \pkg{citation-style-language} package is as follows.

\begin{LaTeXdemo}
  \documentclass{...}
  \usepackage{citation-style-language}
  \cslsetup{
    style = ...,
    ...
  }
  \addbibresource{bibfile.json}
  \begin{document}
  \cite{...}
  ...
  \printbibliography
  \end{document}
\end{LaTeXdemo}

The procedure to compile the document is different across engines.

\paragraph{LuaTeX}

The CSL processor is written in Lua and it can be run directly in LuaTeX
without the need of running external programs.
For LuaTeX, the compiling procedure is simply running \file{latex} twice,
which is the same as documents with cross references.

\paragraph{Other engines}

For engines other than LuaTeX, the \file{citeproc-lua} executable is required
to run on the \file{.aux} file to generate the citations and bibliography.
The general procedure is similar to the traditional BibTeX workflow.
\begin{enumerate}
  \item Run \file{latex} on \file{example.tex}.
  \item Run \file{citeproc-lua} on \file{example.aux}.
    The engine reads the \file{.csl} style, CSL locale files, and
    \file{.bib} database and then writes the processed citations and
      bibliography to \file{example.bbl}.
  \item Run \file{latex} on \file{example.tex}.
    The \file{.bbl} file is loaded and all the citations and bibliography
    are printed.
\end{enumerate}



\section{Global settings}

\begin{function}{\cslsetup}
  \begin{syntax}
    \cs{cslsetup}\marg{options}
  \end{syntax}
\end{function}

Package options may be set when the package is loaded or at any later stage
with the \cs{cslsetup} command.
These two methods are equivalent.
\begin{LaTeXdemo}
  \usepackage[style = apa]{citation-style-langugage}
  % OR
  \usepackage{citation-style-langugage}
  \cslsetup{style = apa}
\end{LaTeXdemo}

\DescribeOption{style}
The \opt{style=}\meta{style-id} option selects the style file
\meta{style-id}\file{.csl} for both citations and bibliography.
The implemented CSL style files are available in the official GitHub
repository\footnote{\url{https://github.com/citation-style-language/styles}}
as well as the Zotero style
repository\footnote{\url{https://www.zotero.org/styles}}.
The user may search and download the \file{.csl} file to the working directory.
The following styles are distributed within the package and
each of them can be directly loaded without downloading.

\begin{description}
  \item[\opt{american-chemical-society}] American Chemical Society
  \item[\opt{american-medical-association}] American Medical Association 11th edition
  \item[\opt{american-political-science-association}] American Political Science Association
  \item[\opt{american-sociological-association}] American Sociological Association 6th edition
  \item[\opt{apa}] American Psychological Association 7th edition
  \item[\opt{chicago-author-date}] Chicago Manual of Style 17th edition (author-date)
  \item[\opt{chicago-fullnote-bibliography}] Chicago Manual of Style 17th edition (full note)
  \item[\opt{chicago-note-bibliography}] Chicago Manual of Style 17th edition (note)
  \item[\opt{elsevier-harvard}] Elsevier - Harvard (with titles)
  \item[\opt{harvard-cite-them-right}] Cite Them Right 11th edition - Harvard
  \item[\opt{ieee}] IEEE
  \item[\opt{modern-humanities-research-association}] Modern Humanities Research Association 3rd edition (note with bibliography)
  \item[\opt{modern-language-association}] Modern Language Association 9th edition
  \item[\opt{nature}] Nature
  \item[\opt{vancouver}] Vancouver
\end{description}

\DescribeOption{locale}
The \opt{locale} option receives an ISO 639-1 two-letter language code
(e.g.,  ``\opt{en}'', ``\opt{zh}''), optionally with a two-letter locale code
(e.g., ``\opt{de-DE}'', ``\opt{de-AT}'').
This option affects sorting of the entries and the output of dates, numbers,
and terms (e.g., ``et al.'').
It may also be set \opt{auto} (default) and the \opt{default-locale} attribute in
the CSL style file will be used.
The locale falls back to ``\opt{en}'' (English) if the attribute is not set.
When \pkg{babel} package is loaded, the selected main language is implicitly set
as the \opt{locale} for \pkg{citation-style-language}.

\DescribeOption{bib-font}
Usually, the list of references is printed in the same font style and size as
the main text.
The \opt{bib-font} option is used to set different formats in the
\env{thebibliography} environment.
It may override the \opt{line-spacing} attribute configured in the CSL style.
For example, to force double-spacing in the bibliography:
\begin{LaTeXdemo}
  \cslsetup{bib-font = \linespread{2}\selectfont}
\end{LaTeXdemo}

\DescribeOption{bib-item-sep}
The vertical space between entries in the bibliography is configured in the
CSL style.
It can be overridden by this \opt{bib-item-sep} option.
It is recommended to set \opt{bib-item-sep} to a stretchable glue rather than
a fixed length to help reducing page breaks in the middle of an entry.
\begin{LaTeXdemo}
  \cslsetup{bib-item-sep = 8 pt plus 4 pt minus 2 pt}
\end{LaTeXdemo}

\DescribeOption{bib-hang}
The \opt{bib-hang} option sets the hanging indentation length which is
usually used for author-date style references.
By default, it is 1 em (with respect to the \opt{bib-font} size if set).



\section{Bibliographic resources}

\begin{function}{\addbibresource}
  \begin{syntax}
    \cs{addbibresource}\oarg{options}\marg{resource}
  \end{syntax}
\end{function}

The \cs{addbibresource} command adds the contents of \meta{resource} into the
bibliographic metadata.
The \meta{resource} may be a CSL-JSON file or the Bib(La)TeX \file{.bib} file.
CSL-JSON \footnote{\url{https://github.com/citation-style-language/schema\#csl-json-schema}}
is the default data model defined by CSL.
Its contents are usually exported from Zotero.
The traditional \file{.bib} file is converted to CSL-JSON internally for
further processing.
The mapping of entry-types and fields between them is detailed in the GitHub wiki
page\footnote{\url{https://github.com/zepinglee/citeproc-lua/wiki/Bib-CSL-mapping}}.
Note that only UTF-8 encoding is supported in the \meta{resource} file.
\begin{LaTeXdemo}
  \addbibresource{data-file.json}
  \addbibresource{bib-file.bib}
\end{LaTeXdemo}


CSL-JSON is the bibliographic metadata format for CSL styles.
Below is an example entry in this format. Each entry must have these two
fields: \texttt{id} and \texttt{type} and the former is the entry keys that
appears in the \cs{cite} command. The accepted entry types and fields are
listed in Appendix \ref{sec:entry-types} and Appendix \ref{sec:entry-fields},
respectively. In most cases, one does not manually write CSL-JSON but exports
from a reference manager like Zotero or Mendeley.

\begin{lstlisting}[style=plain,gobble=2,frame=single]
  {
    "id": "ITEM-1",
    "type": "book",
    "author": [
      {
        "family": "D'Arcus",
        "given": "Bruce"
      }
    ],
    "issued": {
      "date-parts": [
        [ 2005, 11, 22 ]
      ]
    },
    "publisher": "Routledge",
    "title": "Boundaries of dissent: Protest and state power in the media age"
  }
\end{lstlisting}

The traditional Bib(La)TeX format is also supported. It's internally converted
to CSL-JSON format. However there are some differences.

In a Bib(La)TeX database, the titles (includding \texttt{booktitle}s) should
be written in title case with propper words enclosed with braces (e.g.,
“\texttt{The \{Feynman\} Lectures on Physics}”). However CSL requires the
titles stored in sentence case (“\texttt{The Feynman lectures on physics}”).
By default the BibTeX-to-CSL translation does this case conversion.
It checks the original case and skips the titles that are already in sentence
case. For example, the lowercase word “lectures” in “The Feynman lectures on
physics” indicates the whole title is in sentence case and citeproc-lua does
not convert it. This behavior can be turned off by specifying
\opt{check-sentence-case = false} in the optional argument of
\cs{addbibresource}.


\section{Citation commands}

\begin{function}{\cite}
  \begin{syntax}
    \cs{cite}\oarg{options}\marg{keys}
  \end{syntax}
\end{function}

The citation command is similar to the one in standard LaTeX except that the
\meta{options} is in key-value style.
\DescribeOption{prefix}
\DescribeOption{suffix}
\DescribeOption{page}
\DescribeOption{figure}
The \meta{options} can be \opt{prefix}, \opt{suffix} or one of locators like
\opt{page} or \opt{figure}.
The full list of supported locators is detailed in Table~\ref{tab:locators}.
An example is as follows.
\begin{LaTeXdemo}
  \cite[prefix = {See }, page = 42]{ITEM-1}
\end{LaTeXdemo}

\begin{table}
  \centering
  \caption{The locators supported in CSL v1.0.2.}
  \label{tab:locators}
  \begin{tabular}{lll}
    \toprule
    \opt{act}       & \opt{folio}     & \opt{section}    \\
    \opt{appendix}  & \opt{issue}     & \opt{sub-verbo}  \\
    \opt{article}   & \opt{line}      & \opt{supplement} \\
    \opt{book}      & \opt{note}      & \opt{table}      \\
    \opt{canon}     & \opt{opus}      & \opt{timestamp}  \\
    \opt{chapter}   & \opt{page}      & \opt{title}      \\
    \opt{column}    & \opt{paragraph} & \opt{verse}      \\
    \opt{elocation} & \opt{part}      & \opt{version}    \\
    \opt{equation}  & \opt{rule}      & \opt{volume}     \\
    \opt{figure}    & \opt{scene}     &                  \\
    \bottomrule
  \end{tabular}
\end{table}

The traditional form \cs{cite}\oarg{prenote}\oarg{postnote}\marg{keys}
introduced in \pkg{natbib} and \pkg{biblatex} is also supported but not
recommended.
If only one optional argument is provided, it is treated as \meta{postnote}.
The \meta{postnote} is used as a page locator if it consists of only digits.

\begin{function}{\parencite,\citep}
  \begin{syntax}
    \cs{parencite}\oarg{options}\marg{keys}
  \end{syntax}
\end{function}

The \cs{parencite} and \cs{citep} command are aliases of \cs{cite}.
They are added for compatibility with \pkg{biblatex} and \pkg{natbib} packages.
If the citation format defined in the CSL style does not have affixes,
these commands in \pkg{citation-style-language} do not enclose the output with
brackets, which is different from other packages.

\begin{function}{\textcite,\citet}
  \begin{syntax}
    \cs{textcite}\oarg{options}\marg{keys}
  \end{syntax}
\end{function}

\DescribeOption{infix}
These commands proceduce narrative in-text citation where the author name is
part of the running text followed by the year in parentheses.
These commands only work with author-date styles.
An extra option \opt{infix} can be given to specify the text inserted between
then author and year parts. For example, “Kesey’s early work (1962)” can be
produced by |\textcite[infix={'s early work}]{ITEM-1}|.
By default the infix is a space.

\begin{function}{\cites}
  \begin{syntax}
    \cs{cites}\oarg{options}\marg{key}...[options]\marg{key}
  \end{syntax}
\end{function}

The \cs{cites} accepts multiple cite items in a single citation.
This command scans greedily for arguments and a following bracket may be
mistakenly recognized as a delimiter.
To prevent this, an explicit \cs{relax} command is required to terminate the
scanning process. The following example illustrates its usage.

\begin{LaTeXdemo}
  \cites[prefix = {See }, page = 6]{key1}[section = 2.3]{key2}\relax [Text]
\end{LaTeXdemo}

\begin{function}{\citeauthor}
  \begin{syntax}
    \cs{citeauthor}\marg{key}
  \end{syntax}
\end{function}

This command prints the author name.
If the orginal citation does not contain the author name (e.g., a numeric
style), an optional |<intext>| element can be suppplied as a sibling to the
|<citation>| and |<bibliography>| elements in the CSL style (see
\href{https://citeproc-js.readthedocs.io/en/latest/running.html#citation-flags-with-processcitationcluster}{citeproc-js's documentation} for details).

\begin{function}{\nocite}
  \begin{syntax}
    \cs{nocite}\marg{keys}
  \end{syntax}
\end{function}

This command produces no output but makes the entries included in the
bibliography, which is the same in standard \LaTeX.
If the special key |*| is given (\cs{notecite\{*\}}), all the entries in the
database are included.


\section{Biliography commands}

\begin{function}{\printbibliography}
  \begin{syntax}
    \cs{printbibliography}\oarg{key = value, \dots}
  \end{syntax}
\end{function}

This command prints the reference list. It takes one optional argument of the
following options in key-value form.

% \DescribeOption{env}

\newcommand\defaultoption[1]{%
  \hfill
  default: \opt{#1}\par
}

\begin{itemize}
  \item \opt{heading = \meta{name}}:
    This option selects the heading \meta{name} defined by \cs{defbibheading}.
    See also \S \ref{sec:bib-heading}. (default: \opt{bibliography})

  \item \opt{title = \meta{text}}:
    This option overrides the title of selected heading.

  \item \opt{type = \meta{csltype}}:
    Prints only entries with CSL type \meta{meta}. Note that the types of
    entries from BibTeX format (.bib) are converted to CSL type.

  \item \opt{nottype = \meta{csltype}}:
    Exclude entries with CSL type \meta{meta}. Note that the types of
    entries from BibTeX format (.bib) are converted to CSL type.
    This option may be used multiple times.

  \item \opt{keyword = \meta{keyword}}:
    Prints only entries whose \opt{keyword} field includes \meta{keyword}.

  \item \opt{notkeyword = \meta{keyword}}:
    Exclude entries with \meta{keyword} in \opt{keyword} field.
    This option may be used multiple times.

  \item \opt{category = \meta{category}}:
    Prints only entries in the category \meta{category} defined by
    \cs{DeclareBibliographyCategory}. See also \S \ref{sec:bib-category}.

  \item \opt{notcategory = \meta{category}}:
    Exclude entries in the category \meta{category}.
    This option may be used multiple times.

\end{itemize}


\subsection{Bibliography headings}

\begin{function}{\defbibheading}
  \begin{syntax}
    \cs{defbibheading}\marg{name}\oarg{title}\marg{code}
  \end{syntax}
\end{function}

This command defines a bibliography heading \meta{name} which can be selected
with \opt{heading} option in \cs{printbibliography}.
The \meta{code} is expected to generate a section heading with commands like
\cs{chapter} or \cs{section}
The optional \meta{title} argument is the default title passed to \meta{code}
as |#1|.
If the \opt{title} option is used in \cs{printbibliography}, its value will
override the \meta{title}.

Here is the the default \opt{bibliography} heading for standard document class
\cls{article}.

\begin{LaTeXdemo}
  \defbibheading{bibliography}[\refname]{%
    \section*{#1}%
    \@mkboth{\MakeUppercase{#1}}{\MakeUppercase{#1}}%
  }
\end{LaTeXdemo}

The following are predefined headings that can be directly used.

\begin{description}
  \item[\opt{bibliography}]
    This is the default heading for \cs{printbibliography}.
    Its actaul definition depends on the document class.
    For classes like \cls{book} and \cls{report} that define \cs{chapter},
    this heading uses \cs{chapter*} and \cs{bibname} as the default heading
    title. For classes without \cs{chapter}, the heading uses \cs{section*} and
    \cs{refname}.

  \item[\opt{subbibliography}]
    Similar to \opt{bibliography} but one sectioning level lower.
    For example, this heading calls \cs{subsection*} instead of the original
    \cs{section*} in \cls{article} class.

  \item[\opt{bibintoc}]
    Similar to \opt{bibliography} above but adds an entry to the table of
    contents.

  \item[\opt{subbibintoc}]
    Similar to \opt{subbibliography} above but adds an entry to the table of
    contents.

  \item[\opt{bibnumbered}]
    This heading uses a numbered sectioning command like \cs{chapter} or
    \cs{section}. It also adds an entry in the table of contents.

  \item[\opt{subbibnumbered}]
    Similar to \opt{bibnumbered} but one sectioning level lower.

  \item[\opt{none}]
    An empty heading.

\end{description}



\subsection{Bibliography categories}

\begin{function}{\DeclareBibliographyCategory}
  \begin{syntax}
    \cs{DeclareBibliographyCategory}\marg{category}
  \end{syntax}
\end{function}

This command declares a new \meta{category} for use with the \opt{category} and
\opt{notcategory} options in \cs{printbibliography}.

\begin{function}{\addtocategory}
  \begin{syntax}
    \cs{addtocategory}\marg{category}\marg{entrykey, \dots}
  \end{syntax}
\end{function}

This command assigns the \meta{entrykey}s to \meta{category}.




% \begin{function}{\cites}
%   \begin{syntax}
%     \cs{cite}\oarg{options}\marg{keys}
%   \end{syntax}
% \end{function}




% \markdownInput{bib-csl-mapping.md}


\section{Compatibility with other packages}

\paragraph{\pkg{babel}}

The main language set by \pkg{babel} is used as the locale for \pkg{citation-style-language}.
In general, \pkg{babel} is supposed to be loaded before \pkg{citation-style-language}.

\paragraph{\pkg{hyperref}}

When \pkg{hyperref} is loaded, the DOIs, PMIDs, and PMCIDs are correctly
rendered as hyperlinks.
But the citations are not linked to the entries in bibliography.

\paragraph{Incompatible packages}

The following packages are not compatible with \pkg{citation-style-language}.
An error will be triggered if any of them is loaded together with \pkg{citation-style-language}.
\begin{itemize}
  \item \pkg{babelbib}
  \item \pkg{backref}
  \item \pkg{biblatex}
  \item \pkg{bibtopic}
  \item \pkg{bibunits}
  \item \pkg{chapterbib}
  \item \pkg{cite}
  \item \pkg{citeref}
  \item \pkg{inlinebib}
  \item \pkg{jurabib}
  \item \pkg{mcite}
  \item \pkg{mciteplus}
  \item \pkg{multibib}
  \item \pkg{natbib}
  \item \pkg{splitbib}
\end{itemize}



\section{To-do list}

The \pkg{citation-style-language} package is in early development stage and
some features may not work as expected.
Bug report are welcome at the GitHub issue tracker
\footnote{\url{https://github.com/zepinglee/citeproc-lua/issues}}.
The following is a list of features to be implemented.
If you need any of them or new features, please post a issue to let me know
so that I can give it a priority.

\begin{itemize}
  \item The \pkg{citeproc-lua} engine has not passed all the fixures from the
    CSL standard test-suite. The skipped fixtures are lists in
    \href{https://github.com/zepinglee/citeproc-lua/blob/main/tests/citeproc-test-skip.txt}{citeproc-test-skip.txt}
    and they need to be handled (though less than 6\% of test-suite).

  \item The Unicode sorting method is provided by \pkg{lua-uca} package and
    CJK scripts are not supported so far.

  \item Citation commands that capitalize the first letter (\cs{Cite},
    \cs{Textcite}, etc.).

  \item Back references: page numbers of the citations appears after the entry
    item in bibliography (even without the hyperref).

  \item \cs{footcite} command.

  \item \cs{cite} in a footnote should work as in-text citations (similar to
    \cs{smartcite}).

  \item CSL-YAML support.

  \item Multiple bibliographies in a document like \pkg{chapterbib} package or
    \texttt{refsection} and \texttt{refsegment} options in \pkg{biblatex}.

  \item Journal abbreviation.

  \item Sentences case conversion: the title and booktitle fields in BibTeX
    database are converted to sentences case if they are in title case.

  \item Distinguish dropping and non-dropping particles in names.
\end{itemize}



\appendix

\section{CSL-JSON entry types}
\label{sec:entry-types}

The following content is taken from
\url{https://docs.citationstyles.org/en/stable/specification.html#appendix-iii-types}.

\begin{description}
  \item[article]
    A self-contained work made widely available but not published in a
    journal or other publication;\\
    Use for preprints, working papers, and similar works posted on a
    platform where some level of persistence or stewardship is expected
    (e.g. arXiv or other preprint repositories, working paper series);\\
    For unpublished works not made widely available or only hosted on
    personal websites, use \texttt{manuscript}

  \item[article-journal]
    An article published in an academic journal

  \item[article-magazine]
    An article published in a non-academic magazine

  \item[article-newspaper]
    An article published in a newspaper

  \item[bill]
    A proposed piece of legislation

  \item[book]
    A book or similar work;\\
    Can be an authored book or an edited collection of self-contained
    chapters;\\
    Can be a physical book or an ebook;\\
    The format for an ebook may be specified using \texttt{medium};\\
    Can be a single-volume work, a multivolume work, or one volume of a
    multivolume work;\\
    If a \texttt{container-title} is present, the item is interpreted as a
    book republished in a collection or anthology;\\
    Also used for whole conference proceedings volumes or exhibition
    catalogs by specifying \texttt{event} and related variables

  \item[broadcast]
    A recorded work broadcast over an electronic medium (e.g. a radio
    broadcast, a television show, a podcast);\\
    The type of broadcast may be specified using \texttt{genre};\\
    If \texttt{container-title} is present, the item is interpreted as an
    episode contained within a larger broadcast series (e.g. an episode in a
    television show or an episode of a podcast)

  \item[chapter]
    A part of a book cited separately from the book as a whole (e.g. a
    chapter in an edited book);\\
    Also used for introductions, forewords, and similar supplemental
    components of a book

  \item[classic]
    A classical or ancient work, sometimes cited using a common abbreviation

  \item[collection]
    An archival collection in a museum or other institution

  \item[dataset]
    A data set or a similar collection of (mostly) raw data

  \item[document]
    A catch-all category for items not belonging to other types;\\
    Use a more specific type when appropriate

  \item[entry]
    An entry in a database, directory, or catalog;\\
    For entries in a dictionary, use \texttt{entry-dictionary};\\
    For entries in an encyclopedia, use \texttt{entry-encyclopedia}

  \item[entry-dictionary]
    An entry in a dictionary

  \item[entry-encyclopedia]
    An entry in an encyclopedia or similar reference work

  \item[event]
    An organized event (e.g., an exhibition or conference);\\
    Use for direct citations to the event, rather than to works contained
    within an event (e.g. a \texttt{presentation} in a conference, a
    \texttt{graphic} in an exhibition) or based on an event (e.g. a
    \texttt{paper-conference} published in a proceedings, an exhibition
    catalog)

  \item[figure]
    A illustration or representation of data, typically as part of a journal
    article or other larger work;\\
    May be in any format (e.g. image, video, audio recording, 3D model);\\
    The format of the item can be specified using \texttt{medium}

  \item[graphic]
    A still visual work;\\
    Can be used for artwork or other works (e.g. journalistic or historical
    photographs);\\
    Can be used for any still visual work (e.g. photographs, drawings,
    paintings, sculptures, clothing);\\
    The format of the item can be specified using \texttt{medium}

  \item[hearing]
    A hearing by a government committee or transcript thereof

  \item[interview]
    An interview of a person;\\
    Also used for a recording or transcript of an interview; \texttt{author}
    is interpreted as the interviewee

  \item[legal\_case]
    A legal case

  \item[legislation]
    A law or resolution enacted by a governing body

  \item[manuscript]
    An unpublished manuscript;\\
    Use for both modern unpublished works and classical manuscripts;\\
    For working papers, preprints, and similar works posted to a repository,
    use \texttt{article}

  \item[map]
    A geographic map

  \item[motion\_picture]
    A video or visual recording;\\
    If a \texttt{container-title} is present, the item is interpreted as a
    part contained within a larger compilation of recordings (e.g. a part of
    a multipart documentary))

  \item[musical\_score]
    The printed score for a piece of music;\\
    For a live performance of the music, use \texttt{performance};\\
    For recordings of the music, use \texttt{song} (for audio recordings) or
    \texttt{motion\_picture} (for video recordings)

  \item[pamphlet]
    A fragment, historical document, or other unusually-published or
    ephemeral work (e.g. a sales brochure)

  \item[paper-conference]
    A paper formally published in conference proceedings;\\
    For papers presented at a conference, but not published in a
    proceedings, use \texttt{speech}

  \item[patent]
    A patent for an invention

  \item[performance]
    A live performance of an artistic work;\\
    For non-artistic presentations, use \texttt{speech};\\
    For recordings of a performance, use \texttt{song} or
    \texttt{motion\_picture}

  \item[periodical]
    A full issue or run of issues in a periodical publication (e.g. a
    special issue of a journal)

  \item[personal\_communication]
    Personal communications between multiple parties;\\
    May be unpublished (e.g. private correspondence between two researchers)
    or collected/published (e.g. a letter published in a collection)

  \item[post]
    A post on a online forum, social media platform, or similar platform;\\
    Also used for comments posted to online items

  \item[post-weblog]
    A blog post

  \item[regulation]
    An administrative order from any level of government

  \item[report]
    A technical report, government report, white paper, brief, or similar
    work distributed by an institution;\\
    Also used for manuals and similar technical documentation (e.g. a
    software, instrument, or test manual);\\
    If a \texttt{container-title} is present, the item is interpreted as a
    chapter contained within a larger report

  \item[review]
    A review of an item other than a book (e.g. a film review, posted peer
    review of an article);\\
    If \texttt{reviewed-title} is absent, \texttt{title} is taken to be the
    title of the reviewed item

  \item[review-book]
    A review of a book;\\
    If \texttt{reviewed-title} is absent, \texttt{title} is taken to be the
    title of the reviewed book

  \item[software]
    A computer program, app, or other piece of software

  \item[song]
    An audio recording;\\
    Can be used for any audio recording (not only music);\\
    If a \texttt{container-title} is present, the item is interpreted as a
    track contained within a larger album or compilation of recordings

  \item[speech]
    A speech or other presentation (e.g. a paper, talk, poster, or symposium
    at a conference);\\
    Use \texttt{genre} to specify the type of presentation;\\
    Use \texttt{event} to indicate the event where the presentation was made
    (e.g. the conference name);\\
    Use \texttt{container-title} if the presentation is part of a larger
    session (e.g. a paper in a symposium);\\
    For papers published in conference proceedings, use
    \texttt{paper-conference};\\
    For artistic performances, use \texttt{performance}

  \item[standard]
    A technical standard or similar set of rules or norms

  \item[thesis]
    A thesis written to satisfy requirements for a degree;\\
    Use \texttt{genre} to specify the type of thesis

  \item[treaty]
    A treaty agreement among political authorities

  \item[webpage]
    A website or page on a website;\\
    Intended for sources which are intrinsically online; use a more specific
    type when appropriate (e.g. \texttt{article-journal},
    \texttt{post-weblog}, \texttt{report}, \texttt{entry});\\
    If a \texttt{container-title} is present, the item is interpreted as a
    page contained within a larger website

\end{description}



\section{CSL-JSON entry fields}
\label{sec:entry-fields}

\subsection{Standard Variables}

The following content is taken from
\url{https://docs.citationstyles.org/en/stable/specification.html#appendix-iv-variables}.

\begin{description}
  \item[abstract]
    Abstract of the item (e.g. the abstract of a journal article)

  \item[annote]
    Short markup, decoration, or annotation to the item (e.g., to indicate
    items included in a review);\\
    For descriptive text (e.g., in an annotated bibliography), use
    \texttt{note} instead

  \item[archive]
    Archive storing the item

  \item[archive\_collection]
    Collection the item is part of within an archive

  \item[archive\_location]
    Storage location within an archive (e.g. a box and folder number)

  \item[archive-place]
    Geographic location of the archive

  \item[authority]
    Issuing or judicial authority (e.g. “USPTO” for a patent, “Fairfax
    Circuit Court” for a legal case)

  \item[call-number]
    Call number (to locate the item in a library)

  \item[citation-key]
    Identifier of the item in the input data file (analogous to BibTeX
    entrykey);\\
    Use this variable to facilitate conversion between word-processor and
    plain-text writing systems;\\
    For an identifer intended as formatted output label for a citation (e.g.
    “Ferr78”), use \texttt{citation-label} instead

  \item[citation-label]
    Label identifying the item in in-text citations of label styles (e.g.
    “Ferr78”);\\
    May be assigned by the CSL processor based on item metadata;\\
    For the identifier of the item in the input data file, use
    \texttt{citation-key} instead

  \item[collection-title]
    Title of the collection holding the item (e.g. the series title for a
    book; the lecture series title for a presentation)

  \item[container-title]
    Title of the container holding the item (e.g. the book title for a book
    chapter, the journal title for a journal article; the album title for a
    recording; the session title for multi-part presentation at a
    conference)

  \item[container-title-short]
    Short/abbreviated form of \texttt{container-title};\\
    Deprecated; use \texttt{variable="container-title"\ form="short"}
    instead

  \item[dimensions]
    Physical (e.g. size) or temporal (e.g. running time) dimensions of the
    item

  \item[division]
    Minor subdivision of a court with a \texttt{jurisdiction} for a legal
    item

  \item[DOI]
    Digital Object Identifier (e.g. “10.1128/AEM.02591-07”)

  \item[event]
    Deprecated legacy variant of \texttt{event-title}

  \item[event-title]
    Name of the event related to the item (e.g. the conference name when
    citing a conference paper; the meeting where presentation was made)

  \item[event-place]
    Geographic location of the event related to the item (e.g. “Amsterdam,
    The Netherlands”)

  \item[genre]
    Type, class, or subtype of the item (e.g. “Doctoral dissertation” for a
    PhD thesis; “NIH Publication” for an NIH technical report);\\
    Do not use for topical descriptions or categories (e.g. “adventure” for
    an adventure movie)

  \item[ISBN]
    International Standard Book Number (e.g. “978-3-8474-1017-1”)

  \item[ISSN]
    International Standard Serial Number

  \item[jurisdiction]
    Geographic scope of relevance (e.g. “US” for a US patent; the court
    hearing a legal case)

  \item[keyword]
    Keyword(s) or tag(s) attached to the item

  \item[language]
    The language of the item;\\
    Should be entered as an ISO 639-1 two-letter language code (e.g. “en”,
    “zh”), optionally with a two-letter locale code (e.g. “de-DE”, “de-AT”)

  \item[license]
    The license information applicable to an item (e.g. the license an
    article or software is released under; the copyright information for an
    item; the classification status of a document)

  \item[medium]
    Description of the item\textquotesingle s format or medium (e.g. “CD”,
    “DVD”, “Album”, etc.)

  \item[note]
    Descriptive text or notes about an item (e.g. in an annotated
    bibliography)

  \item[original-publisher]
    Original publisher, for items that have been republished by a different
    publisher

  \item[original-publisher-place]
    Geographic location of the original publisher (e.g. “London, UK”)

  \item[original-title]
    Title of the original version
    % (e.g. “Война и мир”, the untranslated Russian title of “War and Peace”)

  \item[part-title]
    Title of the specific part of an item being cited

  \item[PMCID]
    PubMed Central reference number

  \item[PMID]
    PubMed reference number

  \item[publisher]
    Publisher

  \item[publisher-place]
    Geographic location of the publisher

  \item[references]
    Resources related to the procedural history of a legal case or
    legislation;\\
    Can also be used to refer to the procedural history of other items (e.g.
    “Conference canceled” for a presentation accepted as a conference that
    was subsequently canceled; details of a retraction or correction notice)

  \item[reviewed-genre]
    Type of the item being reviewed by the current item (e.g. book, film)

  \item[reviewed-title]
    Title of the item reviewed by the current item

  \item[scale]
    Scale of e.g. a map or model

  \item[source]
    Source from whence the item originates (e.g. a library catalog or
    database)

  \item[status]
    Publication status of the item (e.g. “forthcoming”; “in press”; “advance
    online publication”; “retracted”)

  \item[title]
    Primary title of the item

  \item[title-short]
    Short/abbreviated form of \texttt{title};\\
    Deprecated; use \texttt{variable="title"\ form="short"} instead

  \item[URL]
    Uniform Resource Locator (e.g.
    "\url{https://aem.asm.org/cgi/content/full/74/9/2766}")

  \item[volume-title]
    Title of the volume of the item or container holding the item;\\
    Also use for titles of periodical special issues, special sections, and
    the like

  \item[year-suffix]
    Disambiguating year suffix in author-date styles (e.g. “a” in “Doe,
    1999a”)

\end{description}


\subsection{Number Variables}

Number variables are a subset of the Standard Variables.

\begin{description}
  \item[chapter-number]
    Chapter number (e.g. chapter number in a book; track number on an album)

  \item[citation-number]
    Index (starting at 1) of the cited reference in the bibliography
    (generated by the CSL processor)

  \item[collection-number]
    Number identifying the collection holding the item (e.g. the series
    number for a book)

  \item[edition]
    (Container) edition holding the item (e.g. “3” when citing a chapter in
    the third edition of a book)

  \item[first-reference-note-number]
    Number of a preceding note containing the first reference to the item;\\
    Assigned by the CSL processor;\\
    Empty in non-note-based styles or when the item hasn\textquotesingle t
    been cited in any preceding notes in a document

  \item[issue]
    Issue number of the item or container holding the item (e.g. “5” when
    citing a journal article from journal volume 2, issue 5);\\
    Use \texttt{volume-title} for the title of the issue, if any

  \item[locator]
    A cite-specific pinpointer within the item (e.g. a page number within a
    book, or a volume in a multi-volume work);\\
    Must be accompanied in the input data by a label indicating the locator
    type (see the Locators term list), which
    determines which term is rendered by \texttt{cs:label} when the
    \texttt{locator} variable is selected.

  \item[number]
    Number identifying the item (e.g. a report number)

  \item[number-of-pages]
    Total number of pages of the cited item

  \item[number-of-volumes]
    Total number of volumes, used when citing multi-volume books and such

  \item[page]
    Range of pages the item (e.g. a journal article) covers in a container
    (e.g. a journal issue)

  \item[page-first]
    First page of the range of pages the item (e.g. a journal article)
    covers in a container (e.g. a journal issue)

  \item[part-number]
    Number of the specific part of the item being cited (e.g. part 2 of a
    journal article);\\
    Use \texttt{part-title} for the title of the part, if any

  \item[printing-number]
    Printing number of the item or container holding the item

  \item[section]
    Section of the item or container holding the item (e.g. “§2.0.1” for a
    law; “politics” for a newspaper article)

  \item[supplement-number]
    Supplement number of the item or container holding the item (e.g. for
    secondary legal items that are regularly updated between editions)

  \item[version]
    Version of the item (e.g. “2.0.9” for a software program)

  \item[volume]
    Volume number of the item (e.g. “2” when citing volume 2 of a book) or
    the container holding the item (e.g. “2” when citing a chapter from
    volume 2 of a book);\\
    Use \texttt{volume-title} for the title of the volume, if any

\end{description}


\subsection{Date Variables}

\begin{description}
  \item[accessed]
    Date the item has been accessed

  \item[available-date]
    Date the item was initially available (e.g. the online publication date
    of a journal article before its formal publication date; the date a
    treaty was made available for signing)

  \item[event-date]
    Date the event related to an item took place

  \item[issued]
    Date the item was issued/published

  \item[original-date]
    Issue date of the original version

  \item[submitted]
    Date the item (e.g. a manuscript) was submitted for publication

\end{description}


\subsection{Name Variables}

\begin{description}
  \item[author]
    Author

  \item[chair]
    The person leading the session containing a presentation (e.g. the
    organizer of the \texttt{container-title} of a \texttt{speech})

  \item[collection-editor]
    Editor of the collection holding the item (e.g. the series editor for a
    book)

  \item[compiler]
    Person compiling or selecting material for an item from the works of
    various persons or bodies (e.g. for an anthology)

  \item[composer]
    Composer (e.g. of a musical score)

  \item[container-author]
    Author of the container holding the item (e.g. the book author for a
    book chapter)

  \item[contributor]
    A minor contributor to the item; typically cited using “with” before the
    name when listed in a bibliography

  \item[curator]
    Curator of an exhibit or collection (e.g. in a museum)

  \item[director]
    Director (e.g. of a film)

  \item[editor]
    Editor

  \item[editorial-director]
    Managing editor (“Directeur de la Publication” in French)

  \item[editor-translator]
    Combined editor and translator of a work;\\
    The citation processory must be automatically generate if
    \texttt{editor} and \texttt{translator} variables are identical;\\
    May also be provided directly in item data

  \item[executive-producer]
    Executive producer (e.g. of a television series)

  \item[guest]
    Guest (e.g. on a TV show or podcast)

  \item[host]
    Host (e.g. of a TV show or podcast)

  \item[illustrator]
    Illustrator (e.g. of a children\textquotesingle s book or graphic novel)

  \item[interviewer]
    Interviewer (e.g. of an interview)

  \item[narrator]
    Narrator (e.g. of an audio book)

  \item[organizer]
    Organizer of an event (e.g. organizer of a workshop or conference)

  \item[original-author]
    The original creator of a work (e.g. the form of the author name listed
    on the original version of a book; the historical author of a work; the
    original songwriter or performer for a musical piece; the original
    developer or programmer for a piece of software; the original author of
    an adapted work such as a book adapted into a screenplay)

  \item[performer]
    Performer of an item (e.g. an actor appearing in a film; a muscian
    performing a piece of music)

  \item[producer]
    Producer (e.g. of a television or radio broadcast)

  \item[recipient]
    Recipient (e.g. of a letter)

  \item[reviewed-author]
    Author of the item reviewed by the current item

  \item[script-writer]
    Writer of a script or screenplay (e.g. of a film)

  \item[series-creator]
    Creator of a series (e.g. of a television series)

  \item[translator]
    Translator

\end{description}




\end{documentation}

\end{document}
